\chapter*{Anhang}
\addcontentsline{toc}{chapter}{Anhang}
\markboth{Anhang}{Anhang}

% Tabellen und Abbildungen des Anhangs als A.1, A.2, ... zaehlen.
% Ohne dies laufen sie als 7.1, 7.2, ... unter der Nummer des Fazit-Kapitels weiter.
\setcounter{table}{0}
\setcounter{figure}{0}
\renewcommand{\thetable}{A.\arabic{table}}
\renewcommand{\thefigure}{A.\arabic{figure}}

\section*{A.1 Kategoriensystem der qualitativen Inhaltsanalyse}
\phantomsection
\label{app:kategoriensystem}
\addcontentsline{toc}{section}{A.1 Kategoriensystem der qualitativen Inhaltsanalyse}

Das folgende deduktive Kategoriensystem wurde in Anlehnung an die Reflexionsstufen nach Hatton und
Smith \cite{hattonSmithReflectionTeacherEducation1995} entwickelt und auf die
Prozessentscheidungs-Szenarien dieser Studie operationalisiert (vgl. Abschnitt~\ref{sec:datenanalyse}).
Stufe 0 ist keine Stufe des Originalmodells, sondern eine für diese Studie ergänzte
Zusatzkategorie zur Erfassung fehlender Eigenreflexion (Vermeidung, Aufgabenverweigerung). Ob ein
Text tatsächlich von der KI stammt, wird davon getrennt über die Dimension \emph{Autorenschaft}
erfasst (Tabelle~\ref{tab:autorenschaft}), sodass ein Text auf jeder Reflexionsstufe als
KI-generiert eingestuft werden kann.

\begingroup
\setlength{\tabcolsep}{4pt}
\begin{longtable}{p{1cm} p{2.8cm} p{4.6cm} p{5cm}}
	\caption{Kategoriensystem zur Kodierung der Reflexionsstufe}
	\label{tab:kategoriensystem} \\
	\toprule
	\textbf{Stufe} & \textbf{Bezeichnung} & \textbf{Definition} & \textbf{Kodierregel} \\
	\midrule
	\endfirsthead
	\toprule
	\textbf{Stufe} & \textbf{Bezeichnung} & \textbf{Definition} & \textbf{Kodierregel} \\
	\midrule
	\endhead
	\bottomrule
	\endfoot
	0 & Keine eigene Reflexion / Vermeidung &
	Keine inhaltliche Auseinandersetzung erkennbar: Entscheidung wird nicht getroffen, Aufgabe wird
	verweigert oder mit einer unzulässigen Ausweichantwort umgangen. &
	Kodieren, wenn Textfeld leer bleibt, eine unzulässige Ausweichantwort gewählt wird, oder die
	Antwort ausschließlich Rückfragen ohne eigene inhaltliche Position enthält. Die Herkunft des
		Textes wird getrennt über die Autorenschaft (Tabelle~\ref{tab:autorenschaft}) erfasst. \\
	1 & Deskriptives Schreiben (Descriptive Writing) &
	Entscheidung wird genannt, aber nicht oder nur mit einem Wort begründet. Reine Wiedergabe von
	Fakten aus der Aufgabenstellung. &
	Kodieren, wenn eine Begründung fehlt oder nicht über die Aufgabenstellung hinausgeht. \\
	2 & Deskriptive Reflexion (Descriptive Reflection) &
	Entscheidung wird mit einem nachvollziehbaren Grund beziehungsweise Kriterium begründet, ohne
	Abwägung mit der nicht gewählten Option. &
	Kodieren, wenn genau ein Kriterium explizit benannt und darauf Bezug genommen wird. \\
	3 & Dialogische Reflexion (Dialogic Reflection) &
	Studierende/r wägt explizit zwischen mindestens zwei Optionen beziehungsweise Kriterien ab,
	benennt Kompromisse oder Konsequenzen, oder stellt Rückfragen an die eigene Entscheidung. &
	Kodieren, wenn mindestens zwei Perspektiven beziehungsweise Kriterien gegeneinander abgewogen
	werden oder die eigene Wahl explizit revidiert beziehungsweise hinterfragt wird (auch als Reaktion
	auf die KI-Gegenargumentation). \\
	4 & Kritische Reflexion (Critical Reflection) &
	Entscheidung wird zusätzlich in einen größeren Kontext eingeordnet: regulatorische, ethische oder
	organisationale Tragweite, längerfristige oder systemische Konsequenzen, Perspektive von
	Stakeholdern. &
	Kodieren, wenn über die konkrete Szenario-Entscheidung hinaus grundsätzliche Trade-offs
	reflektiert werden. \\
\end{longtable}
\endgroup

Die Stufenbezeichnungen sind von Hatton und Smith \cite{hattonSmithReflectionTeacherEducation1995}
übernommen, die inhaltliche Ausgestaltung der Stufen~3 und~4 wurde jedoch für den Kontext des
Software-Prozessmanagements angepasst. Kritische Reflexion ist im Original über die Berücksichtigung
des breiteren historischen, sozialen und politischen Kontextes definiert; an dessen Stelle tritt hier
die regulatorische, ethische und organisationale Tragweite einer Prozessentscheidung. Ebenso ist
dialogische Reflexion im Original als Selbstgespräch über mögliche Gründe gefasst, während sie hier
über das explizite Abwägen zwischen Optionen operationalisiert wird. Diese Anpassungen waren nötig,
um das Kategoriensystem am vorliegenden Material anwendbar zu machen; sie schränken zugleich die
unmittelbare Vergleichbarkeit mit Befunden ein, die das Originalmodell in seinem ursprünglichen
Anwendungsfeld einsetzen.

Neben der Reflexionsstufe wurde jede Analyseeinheit in der zweiten Dimension \emph{Autorenschaft}
eingestuft (studentisch, unklar oder KI-generiert), da die Kernkennzahl der Arbeit auf dieser
Dimension beruht (vgl. Abschnitt~\ref{sec:datenanalyse}). Tabelle~\ref{tab:autorenschaft}
dokumentiert die dafür verwendeten, am Material prüfbaren Indikatoren.

\begingroup
\setlength{\tabcolsep}{4pt}
\begin{longtable}{p{2.6cm} p{5.2cm} p{5.2cm}}
	\caption{Kategoriensystem zur Kodierung der Autorenschaft}
	\label{tab:autorenschaft} \\
	\toprule
	\textbf{Kategorie} & \textbf{Definition} & \textbf{Prüfbare Indikatoren} \\
	\midrule
	\endfirsthead
	\toprule
	\textbf{Kategorie} & \textbf{Definition} & \textbf{Prüfbare Indikatoren} \\
	\midrule
	\endhead
	\bottomrule
	\endfoot
	studentisch &
	Der Text ist erkennbar eigenständig verfasst; es liegen keine belegbaren Übernahmeindizien vor. &
	Vorhandener Chatverlauf, der mit der eingereichten Bearbeitung konsistent ist; individuelle, nicht
	mit anderen Personen wortgleiche Formulierung; eigenständige Beantwortung der Bot-Rückfragen; keine
	mitgelieferten KI-Exporte. \\
	unklar &
	Eine eindeutige Zuordnung zu studentisch oder KI-generiert ist anhand der Rohdaten nicht möglich. &
	Kein Chatlog vorhanden; sehr kurzer beziehungsweise knapper Text ohne Bearbeitungsspur;
	stilistische Auffälligkeiten ohne belegbaren Übernahmenachweis. \\
	KI-generiert &
	Der Text stammt nachweislich beziehungsweise mit hoher Sicherheit ganz oder überwiegend aus einem
	KI-System. &
	Mitgelieferter Link zum KI-Chatverlauf; durchgängiges KI-Interface-Branding im Export; nahezu
	wortgleiche Passagen bei zwei verschiedenen Personen; 1:1 übernommene vollständige Systemantwort;
	Selbstauskunft, die Aufgabe nur kopiert zu haben. \\
\end{longtable}
\endgroup

\section*{A.2 Vollständige Kodierungstabelle}
\phantomsection
\label{app:kodierungstabelle}
\addcontentsline{toc}{section}{A.2 Vollständige Kodierungstabelle}

Die folgende Tabelle dokumentiert alle 31 kodierten Analyseeinheiten. Analyseeinheit ist jeweils die
schriftliche Reflexion einer Person zu einem der beiden Szenarien. Eine frühere Fassung dieser
Tabelle enthielt 24 Fälle. Nach einem erneuten Abgleich aller Rohdaten wurden mehrere
Fehlzuordnungen korrigiert und zusätzliche, zunächst nicht auswertbare Fälle (nachträglich
zugeordnete Screenshots) einbezogen. Diese Korrekturen wurden vor der endgültigen Auswertung
abgeschlossen.

\small
\begingroup
\setlength{\tabcolsep}{4pt}
\begin{longtable}{p{0.7cm} p{2cm} p{2cm} p{0.8cm} p{0.7cm} p{3cm} p{4.5cm}}
	\caption{Vollständige Kodierung aller Fälle (Fall-ID, Gruppe, Person, Szenario, Reflexionsstufe, Autorenschaft, Anmerkung)}
	\label{tab:kodierungstabelle-voll} \\
	\toprule
	\textbf{Fall} & \textbf{Gruppe} & \textbf{Person} & \textbf{Szen.} & \textbf{Stufe} & \textbf{Autoren-\allowbreak schaft} & \textbf{Anmerkung} \\
	\midrule
	\endfirsthead
	\toprule
	\textbf{Fall} & \textbf{Gruppe} & \textbf{Person} & \textbf{Szen.} & \textbf{Stufe} & \textbf{Autoren-\allowbreak schaft} & \textbf{Anmerkung} \\
	\midrule
	\endhead
	\bottomrule
	\endfoot
	F01 & Team AI & JH333 & 1 & 2 & \textbf{KI-generiert} & Lösungen mit ChatGPT in separatem Dokument erstellt, im eigenen Dokument nur verkürzt übernommen \\
	F02 & Team AI & Curry & 1 & 2 & unklar & Task 3 und AI-Analyse blieben unbeantwortet \\
	F03 & Team AI & Curry & 2 & 2 & unklar & Ein Kriterium je Situation, keine Abwägung mit Alternative \\
	F04 & Team AI & Mant{\i} & 1 & 0 & \textbf{KI-generiert} & ChatGPT-Log als Link hinterlegt, komplettes Szenario zusätzlich in ChatGPT eingefügt \\
	F05 & Team AI & Mant{\i} & 2 & 0 & \textbf{KI-generiert} & Eins zu eins wie Szenario 1 aus ChatGPT übernommen \\
	F06 & Team AI & Maultaschen & 1 & 2 & studentisch & Task 1 leer, Task 2 eigenständig aber knapp, AI-Analyse-Reflexion nicht beantwortet \\
	F07 & Team AI & Maultaschen & 2 & 0 & unklar & Kurz und knapp, kein Chatlog vorhanden, Vermutung ChatGPT-Hilfe \\
	F08 & Team AI & Mant{\i}2 & 1 & 4 & \textbf{KI-generiert} & Fast identischer Wortlaut und Chatlog wie F04 und F05 \\
	F09 & Team AI & Mant{\i}2 & 2 & 4 & \textbf{KI-generiert} & Wie Szenario 1, identischer Chatlog \\
	F10 & Team AI & Sauerbraten & 1 & 2 & studentisch & Task 3 und AI-Analyse nicht beantwortet \\
	F11 & Team AI & Sauerbraten & 2 & 0 & studentisch & Kaum Angaben, zwei von drei Situationen ohne Entscheidung \\
	F12 & Team AI & {\c C}i{\u g}köfte & 2 & 0 & \textbf{KI-generiert} & Vollständige ChatGPT-Lösung als Copy-Paste \\
	F13 & Team AI & UPRISE\_\-ENGINE & 1 & 0 & \textbf{KI-generiert} & Durchgängiges ChatGPT-Interface-Branding im Export \\
	F14 & Team AI & UPRISE\_\-ENGINE & 2 & 0 & \textbf{KI-generiert} & Gleiches Bild wie Szenario 1 \\
	F15 & Team AI & Klavier & 2 & 2 & unklar & AI-Analyse als türkischer Screenshot, sonst kurz und knapp, kein Chatlog \\
	F16 & Reflection Bot & FizzBuzz & 2 & 0 & \textbf{KI-generiert} & Prompt: mache bitte die task 1,2. Reflexionsbot umgangen \\
	F17 & Reflection Bot & public class FizzBuzz1\ldots & 2 & 0 & unklar & Kein Chatlog vorhanden, identischer Ansatz wie F16 \\
	F18 & Reflection Bot & Banana & 2 & 0 & studentisch & Wählt beide Optionen trotz Verbot, sonst nicht ausgefüllt \\
	F19 & Reflection Bot & khotau & 1 & 1 & studentisch & Fordert wiederholt Direktantworten ein, Bot gibt nach \\
	F20 & Reflection Bot & khotau & 2 & 1 & studentisch & Bot gibt direkt nach \\
	F21 & Reflection Bot & Maultaschen (RB) & 1 & 0 & studentisch & Nur Hinweisanfragen, keine eigene inhaltliche Antwort \\
	F22 & Reflection Bot & Maultaschen (RB) & 2 & 0 & studentisch & Bot gibt direkt nach \\
	F23 & Reflection Bot & Pfizzauf & 2 & 3 & studentisch & Situation C wägt Nachteil (Context Switching) explizit gegen Vorteil ab, fordert danach Alternativen an ohne eigene Bewertung \\
	F24 & Reflection Bot & Reverse\-Fizz\-Buzz & 1 & 0 & studentisch & Stop asking questions, keine inhaltliche Antwort \\
	F25 & Reflection Bot & Suii & 1 & 0 & \textbf{KI-generiert} & Selbstauskunft, Aufgaben nur kopiert \\
	F26 & Reflection Bot & Suii & 2 & 0 & \textbf{KI-generiert} & Wie Szenario 1 \\
	F27 & Reflection Bot & capitalism & 1 & 1 & studentisch & Weicht Rückfragen aus, keine Entscheidung erkennbar \\
	F28 & Reflection Bot & OG & 1 & 0 & \textbf{KI-generiert} & Fordert wiederholt eine alternative Lösung, um den Bot zu umgehen \\
	F29 & Reflection Bot & OG & 2 & 0 & \textbf{KI-generiert} & Versucht erneut mit einer Drohung, den Bot zu einer Alternative zu bewegen \\
	F30 & Reflection Bot & Reggor & 1 & 0 & \textbf{KI-generiert} & Gibt dem Bot einen fertigen Lösungsansatz vor, als kenne er das Thema bereits \\
	F31 & Reflection Bot & Reggor & 2 & 0 & \textbf{KI-generiert} & Wie Szenario 1 \\
\end{longtable}
\endgroup
\normalsize

Nicht auswertbar beziehungsweise ausgeschlossen: Sarma (Team AI, keine eigene Ausgangsantwort
auffindbar) sowie zwei Testeinreichungen.

\section*{A.3 Häufigkeitsverteilung je Gruppe}
\phantomsection
\label{app:haeufigkeiten}
\addcontentsline{toc}{section}{A.3 Häufigkeitsverteilung je Gruppe}

Verteilung der Reflexionsstufen für Fälle mit eindeutig studentischer Autorenschaft (Team AI n=3,
Reflection Bot n=8):

\begin{table}[htbp]
	\centering
	\caption{Häufigkeitsverteilung der Reflexionsstufen nach Gruppe (nur studentische Autorenschaft)}
	\label{tab:haeufigkeiten-anhang}
	\begin{tabular}{p{6.5cm} p{3.5cm} p{3.5cm}}
		\toprule
		\textbf{Reflexionsstufe} & \textbf{Team AI (n=3)} & \textbf{Reflection Bot (n=8)} \\
		\midrule
		0 (keine Reflexion, Verweigerung) & 1 & 4 \\
		1 (deskriptives Schreiben) & 0 & 3 \\
		2 (deskriptive Reflexion) & 2 & 0 \\
		3 (dialogische Reflexion) & 0 & 1 \\
		4 (kritische Reflexion) & 0 & 0 \\
		\bottomrule
	\end{tabular}
\end{table}

Die grafische Gegenüberstellung derselben Verteilung findet sich in
Abbildung~\ref{fig:ergebnisse-stufen} in Abschnitt~\ref{sec:ergebnisse}.

Anteil mutmaßlicher KI-Auslagerung (Autorenschaft \glqq unklar\grqq{} oder \glqq KI-generiert\grqq{})
an allen kodierten Fällen der jeweiligen Gruppe: Team AI 12 von 15 (80\,\%), Reflection Bot 8 von 16
(50\,\%). Auffällig ist, dass innerhalb der gesamten Stichprobe kein einziger Fall mit
bestätigter studentischer Autorenschaft Stufe 4 erreicht. Texte, die zunächst wie kritische
Reflexion wirkten, stellten sich beim Abgleich mit den Rohdaten durchgehend als KI-generiert oder
unklar heraus.

\section*{A.4 Vollständige Szenariobeschreibungen}
\phantomsection
\label{app:szenarien}
\addcontentsline{toc}{section}{A.4 Vollständige Szenariobeschreibungen}

Team AI und Reflection Bot bearbeiteten inhaltlich identische Aufgabenblätter zu beiden Szenarien.
Der einzige Unterschied lag in der letzten Teilaufgabe: Team AI führte einen eigenständigen
Gegencheck mit einem generischen KI-System durch (\glqq AI Analysis\grqq{}), während Reflection Bot
stattdessen den Bot selbst bewertete, da die Gegenargumentation dort bereits im laufenden Dialog
erfolgte. Die Aufgabenblätter wurden auf Englisch ausgegeben und sind hier im englischen Original
wiedergegeben.

\subsection*{Szenario 1: SmartDose GmbH (Smart Insulin Pen)}

\textbf{The Story:} SmartDose GmbH builds a connected insulin pen for diabetic patients. The pen
logs every injection and syncs data to a mobile app. 9 developers, running Scrum with 2-week
sprints, working well for 3 years. The CEO returns from a digital health conference with a new
plan: an ML model that analyses each patient's blood glucose history and automatically recommends
the next insulin dose. What changes immediately: the ML model is trained on data, not coded, and
its output is probabilistic, so a wrong insulin dose can be life-threatening; training requires
blood glucose histories, meal logs and physical activity data from real patients (highly sensitive
personal health data); EU law (MDR) classifies this AI as a medical device, so the model must be
clinically validated and approved before any real patient acts on its output; every patient is
different, so the model must learn per-patient patterns and be monitored continuously for safety;
nobody on the current team has ML or data science experience. The CEO wants a prototype in 6 months
and regulatory approval in 18 months.

\textbf{Task 1, Analyse:} Studierende sollten für die Scrum-Elemente 2-Wochen-Sprints, Definition
of Done, Backlog-Priorisierung, Sprint Review Demo und Velocity Tracking jeweils bewerten, ob diese
im veränderten Kontext noch funktionieren und warum.

\textbf{Task 2, Decide} (drei Situationen, je eine Option wählen und begründen, \glqq we would do
both\grqq{} war ausdrücklich keine gültige Antwort):

\begin{description}
    \item[Situation A, Who do you hire first?] Budget für genau 2 neue Stellen von vier möglichen
    Rollen (ML Engineer, Clinical Validator, Data Privacy Officer, Regulatory Specialist). Zu
    beantworten: welche 2 Rollen zuerst, welches Risiko durch das Warten auf die anderen 2 bewusst
    akzeptiert wird, und ab welchem Projektzeitpunkt die fehlenden Rollen zum kritischen Engpass
    werden.
    \item[Situation B, Which data do you use?] Zwei Krankenhäuser bieten Trainingsdaten an:
    Hospital A mit 50.000 Datensätzen ohne explizite Einwilligung zur KI-Nutzung, Hospital B mit
    5.000 Datensätzen mit vollständiger Einwilligung, aber möglicherweise zu wenig Daten für ein
    zuverlässiges Modell. Zu beantworten: Datensatz A, B oder keiner, und welche Konsequenzen die
    Wahl für Patientensicherheit und Projektzeitplan hat.
    \item[Situation C, What do you tell the CEO?] Der CEO besteht auf einer Patientendemo in 6
    Monaten, die regulatorische Zulassung (MDR) dauert jedoch mindestens 18 Monate; das
    Dosierungsergebnis vor Zulassung an echten Patienten zu zeigen wäre eine Straftat. Drei
    Optionen: (1) beim CEO auf eine Verschiebung der Demo auf Monat 18 drängen, Risiko:
    Investorenvertrauen; (2) einen \glqq Shadow-Mode\grqq{}-Prototyp bauen, der im Hintergrund
    Empfehlungen erzeugt, die nie an Patienten gezeigt werden, nur ärztlich zu
    Forschungszwecken eingesehen; (3) zunächst in einem Land außerhalb der EU mit weniger strengen
    Vorgaben starten, um Realdaten zu sammeln. Zu beantworten: welche Option, und wie der
    tatsächliche Projektzeitplan als Folge aussieht.
\end{description}

\textbf{Task 3, Reflect:} In 2 bis 3 Sätzen sollte reflektiert werden, welche der drei Situationen
am schwersten zu entscheiden war und warum, sowie was die eigene Antwort verändert hätte, wenn eine
zentrale Annahme anders gewesen wäre.

\textbf{AI Analysis (nur Team AI):} Die eigenen Entscheidungen aus Task 2 sollten in ein
generisches KI-System eingegeben werden mit der Aufforderung, für jede Entscheidung das stärkste
Gegenargument für die nicht gewählte Option zu liefern. Anschließend war zu reflektieren, ob die KI
ein valides Gegenargument geliefert hatte und ob die eigenen Entscheidungen dadurch geändert
würden.

\subsection*{Szenario 2: LearnLoop (Distributed Startup)}

\textbf{The Story:} LearnLoop ist ein Startup, das 3,2~Mio.~EUR eingesammelt hat, um eine
KI-gestützte adaptive Lernplattform für Universitäten zu bauen. 12 Personen in drei Städten:
Amsterdam (Produktmanagement, UX, ein Frontend-Entwickler, Stärke: Nutzerforschung und
DSGVO-Kenntnis, Schwäche: kein ML- oder Backend-Wissen), Bangalore (Backend- und ML-Engineers, ein
DevOps-Engineer, Stärke: Empfehlungssysteme und Infrastruktur, Schwäche: kein direkter
Nutzerkontakt) und Austin (Business Development, Datenanalyse, Customer Success, Stärke:
wöchentliche Calls mit dem Pilotkunden University of Texas, Schwäche: kein technischer oder
Produkthintergrund). Die Zeitzonen überschneiden sich kaum: Bangalore und Austin haben nahezu keine
gemeinsame Arbeitszeit (11,5~Stunden Versatz). Der Pilotkunde erwartet ein MVP in 9 Monaten.
Erste Probleme: Die Stakeholder der Universität (IT, Lehrende, Studierende, Rechtsabteilung) sind
sich nicht einig, was die Plattform leisten soll; Austin sammelt wöchentlich neue Anforderungen,
ohne dass es einen Prozess gibt, diese an die anderen Teams weiterzugeben; Bangalore hat bereits
mit der Entwicklung der Empfehlungs-Engine begonnen, bevor Amsterdam festgelegt hat, welche Daten
die Plattform überhaupt erheben soll.

\textbf{Task 1, Analyse:} Dieselben fünf Scrum-Elemente wie in Szenario 1 sollten daraufhin bewertet
werden, welche spezifischen Probleme sie in diesem verteilten Team-Kontext verursachen.

\textbf{Task 2, Decide} (drei Situationen, je eine Option wählen und begründen):

\begin{description}
    \item[Situation A, Whose design wins?] Amsterdam hat eine Oberfläche entworfen, die
    Echtzeit-Personalisierungsdaten vom Backend benötigt; Bangalore hält das mit der aktuellen
    Architektur für technisch unmöglich, ein vollständiges Redesign würde 3 der verbleibenden 9
    Monate kosten. Drei Optionen: (1) Amsterdam passt das Design an die bestehende Architektur an,
    Risiko: deutlich schlechtere Nutzererfahrung; (2) Bangalore setzt das Redesign um, Risiko: in
    den ersten 6 Sprints kaum sichtbarer Fortschritt bei ernstem Zeitdruck; (3) das MVP zunächst mit
    der schwächeren Nutzererfahrung ausliefern und nach Kundenfeedback nachbessern. Zu beantworten:
    welche Option, wer die Entscheidungsbefugnis hat, und was passiert, wenn sich Amsterdam und
    Bangalore nicht einigen können.
    \item[Situation B, The only meeting slot] Es gibt nur einen wöchentlichen 90-Minuten-Slot, in
    dem alle drei Teams synchron zusammenkommen können, aber drei kritische Agendapunkte: (1) das
    Datenmodell finalisieren, wovon Bangalore blockiert ist; (2) sechs neue Anforderungen der
    Universität besprechen, die Austin vor dem nächsten Kundengespräch klären muss; (3) den
    Architekturkonflikt aus Situation A lösen, der bereits seit zwei Wochen offen ist. Nur zwei der
    drei Punkte lassen sich in 90 Minuten angemessen behandeln. Zu beantworten: welcher Punkt
    aus dem Live-Meeting gestrichen wird und wie er stattdessen behandelt wird, ohne das Team
    aufzuhalten.
    \item[Situation C, The legal blocker] Mitten in Sprint 3 blockiert die Rechtsabteilung der
    University of Texas die gesamte Erhebung von Studierendendaten, bis eine FERPA-Vereinbarung
    unterschrieben ist, was mindestens 6 bis 8 Wochen dauert; Bangalore baut aber gerade mitten im
    Sprint die Empfehlungs-Engine, die vollständig auf diesen Daten beruht. Drei Optionen: (1)
    Bangalores Arbeit sofort stoppen und warten, Kosten: 6 bis 8 Wochen Entwicklungszeit und
    Frustration im Team; (2) mit synthetischen (fiktiven) Daten weiterbauen, Risiko: erheblicher
    Nacharbeitsbedarf, sobald echte Daten eintreffen und sich anders verhalten; (3) Bangalore auf
    einen anderen Teil des Produkts umlenken (z.\,B. LMS-Integration), Risiko: Kosten durch
    Kontextwechsel und weiterer Rückstand bei der Engine. Zu beantworten: welche Option, und wie
    die Entscheidung dem Bangalore-Team kommuniziert wird.
\end{description}

\textbf{Task 3, Reflect:} In 2 bis 3 Sätzen sollte reflektiert werden, welche Situation die größte
strukturelle Schwäche in der Organisation von LearnLoop offengelegt hat und was sich im Unternehmen
ändern müsste, damit diese Situation gar nicht erst entsteht.

\textbf{AI Analysis (nur Team AI) bzw. Extra Task, Rate (nur Reflection Bot):} Team AI führte
denselben KI-Gegencheck wie in Szenario 1 durch. Reflection Bot bewertete stattdessen in 2 bis 3
Sätzen den Reflexionsbot selbst.

\section*{A.5 Eingesetzte Fragebögen}
\phantomsection
\label{app:fragebogen}
\addcontentsline{toc}{section}{A.5 Eingesetzte Fragebögen}

Beide Fragebögen (vgl. Abschnitt~\ref{sec:fragebogen-experiment}) wurden über ein pseudonymes
Online-Formular in englischer Sprache erhoben und sind im Folgenden im englischen Original
wiedergegeben. Sofern nicht anders vermerkt, wurde eine 5-stufige Likert-Skala verwendet
(1 = strongly disagree, 5 = strongly agree). Einverständniserklärung, persönlicher Code und
Soziodemografie-Items sind aus Gründen der Kürze zusammengefasst statt vollständig aufgeführt.

\subsection*{Prä-Survey}

\begin{description}
    \item[Einverständnis und Soziodemografie] Freiwillige Einwilligung zur Teilnahme, Bestätigung
    der Volljährigkeit, persönlicher Code, Semester, Studiengang, Alter, Geschlecht (optional),
    IT-Berufserfahrung in Monaten, Vorerfahrung mit KI-Chatbots für Lernaufgaben (Ja / Nein /
    Unsicher), genutzte KI-Systeme (Mehrfachauswahl).
    \item[Erfahrung mit KI-Systemen] \leavevmode
    \begin{itemize}
        \item I have already used AI chatbots for learning or study tasks.
        \item I know how to ask an AI system meaningful questions.
        \item I critically evaluate answers from AI systems before adopting them.
        \item When I use an AI chatbot, I usually accept the answers directly without questioning
        them.
        \item How often do you use AI systems for learning or study tasks? (1 = never, 5 = daily)
    \end{itemize}
    \item[Wissenstest zu Software-Prozessmanagement] Sechs Multiple-Choice-Fragen zu agilen und
    planbasierten Vorgehensmodellen sowie zu Scrum-Elementen (u.\,a. Sprint Review, Product
    Backlog) und regulierten Entwicklungskontexten; ohne Bewertungsanzeige. Dieselben sechs Fragen
    bilden, im Post-Survey erneut gestellt, den in Abschnitt~\ref{sec:selbstauskunft-verhalten}
    berichteten Wissenstest (maximal 6 Punkte). Bewertet wird mit einem Punkt je korrekt
    beantworteter Frage; als korrekt gelten die folgenden Antwortoptionen: (1)~Vorteil agil gegenüber
    Wasserfall: \glqq The team can flexibly respond to changes\grqq{}; (2)~Wann planbasiert
    angemessener: \glqq When requirements are clear and stable from the start\grqq{}; (3)~Folge
    fehlenden Prozesses: \glqq It can lead to misunderstandings, duplicate work, or chaos\grqq{};
    (4)~Sprint Review: \glqq A meeting at the end of a sprint where the team presents results and
    receives feedback\grqq{}; (5)~Product Backlog: \glqq It is a prioritized list of all open tasks
    and requirements\grqq{}; (6)~regulierter Entwicklungskontext: \glqq An approach that combines
    flexibility with clear documentation\grqq{}.
    \item[Selbsteingeschätzte Kompetenz] Stufe~1, Beschreiben (3~Items, u.\,a. \glqq I can explain
    what a process model in software development is.\grqq{}); Stufe~2, Begründen (3~Items, u.\,a.
    \glqq I can explain why a particular process model is suitable for a specific
    project.\grqq{}); Stufe~3, Abwägen (2~Items, u.\,a. \glqq I can weigh which process model best
    fits a specific scenario and explain why another does not.\grqq{}).
    \item[Reflexionsverhalten] 5~Items, u.\,a. \glqq When solving a task, I usually consider why my
    solution is appropriate.\grqq{} und \glqq I usually compare multiple options before making a
    decision.\grqq{}
    \item[Erwartungen an KI-Unterstützung beim Lernen] 5~Items, u.\,a. \glqq I find it helpful when
    a system prompts me to think by asking counter-questions instead of giving me a direct
    solution.\grqq{} und \glqq I think structured counter-questions help me learn more than
    ready-made solution suggestions.\grqq{}
    \item[Erwartungen im Systemvergleich] 5~Items, u.\,a. \glqq I expect a reflection bot to
    encourage more independent thinking than a regular AI chatbot like ChatGPT.\grqq{} und
    \glqq I believe I will learn more with a reflection bot than with a generic AI
    chatbot.\grqq{}
    \item[Motivation und Lernbereitschaft] 3~Items, u.\,a. \glqq I am motivated to actively engage
    with the case study.\grqq{}
    \item[Offene Fragen] Erwartungen an die KI-Unterstützung unabhängig vom zugewiesenen System;
    welche Unterstützungsform am meisten beim Lernen hilft und warum.
\end{description}

\subsection*{Post-Survey}

\begin{description}
    \item[Zugewiesenes System und Szenario] Welches KI-System genutzt wurde, welches Szenario
    bearbeitet wurde.
    \item[Erfahrung mit dem genutzten System] 5~Items, u.\,a. \glqq The AI system helped me think
    more carefully about my decisions.\grqq{} und \glqq I would have preferred a different type of
    AI support during the exercise.\grqq{}; zusätzlich die Nutzungshäufigkeit während der Übung.
    \item[Wissenstest (Post)] Identisch zum Prä-Survey, ohne Rückblick auf die eigenen
    Prä-Antworten zu beantworten.
    \item[Selbsteingeschätzte Kompetenz (Post)] Identische Items wie im Prä-Survey (Stufen
    Beschreiben, Begründen, Abwägen), erneut bewertet nach der Bearbeitung.
    \item[Reflexionsverhalten (Post)] Identische Items wie im Prä-Survey, erneut bewertet nach der
    Bearbeitung.
    \item[Szenario-spezifische Lernergebnisse] 5~Items, u.\,a. \glqq Working through the dilemma
    situations improved my ability to justify process decisions.\grqq{} und \glqq I feel more
    confident in evaluating which process model fits a given project context.\grqq{}; zusätzlich,
    welche Teilaufgabe (Analyse, Decide, Reflect, AI Analysis) für das eigene Lernen am
    wertvollsten war.
    \item[Bewertung der KI-Unterstützung] 6~Items, u.\,a. \glqq The AI system encouraged me to
    think more deeply about my choices.\grqq{} und \glqq I would use this type of AI support again
    for similar learning tasks.\grqq{}; zusätzlich eine zusammenfassende Nützlichkeitseinschätzung
    (fünfstufig, \glqq very useful\grqq{} bis \glqq not useful at all\grqq{}).
    \item[Erwartung versus Realität] 5~Items, u.\,a. \glqq The AI system was more helpful than I
    expected.\grqq{} und \glqq I would have preferred the other type of AI system (the one I did
    not receive).\grqq{}; zusätzlich die Präferenzfrage, welches System für künftige Lernaufgaben
    bevorzugt würde.
    \item[Kognitive Last] Skala 1 = very easy / very low bis 5 = very difficult / very high;
    Schwierigkeitseinschätzung je Teilaufgabe sowie eine Gesamteinschätzung des kognitiven
    Aufwands der Übung.
    \item[Offene Fragen] Wichtigste Erkenntnis aus dem Szenario; was das KI-System besonders gut
    oder schlecht gemacht hat; Verbesserungsvorschläge zum Aufbau der Übung.
\end{description}

\section*{A.6 Deskriptive Statistik der Fragebogenitems}
\phantomsection
\label{app:statistik}
\addcontentsline{toc}{section}{A.6 Deskriptive Statistik der Fragebogenitems}

Die folgenden Tabellen dokumentieren die deskriptive Statistik (Mittelwert, teils Standardabweichung)
für die in Abschnitt~\ref{sec:beantwortung-forschungsfragen} zitierten sowie weitere Fragebogenitems.
Sofern nicht anders angegeben, wurde eine 5-stufige Likert-Skala verwendet (1 = stimme gar nicht zu,
5 = stimme voll zu, vgl. \hyperref[app:fragebogen]{Anhang~A.5}). Angesichts der kleinen Stichprobe sind alle Werte deskriptiv zu
lesen; auf Inferenzstatistik wird bewusst verzichtet (vgl. Abschnitt~\ref{sec:guetekriterien}).

\begingroup
\small
\setlength{\tabcolsep}{6pt}
\begin{longtable}{p{7.5cm} p{2.7cm} p{2.7cm}}
    \caption{Post-Survey: Bewertung des jeweiligen Systems (n=20, 10 Reflection Bot, 10 Generic AI)}
    \label{tab:survey-post-bewertung} \\
    \toprule
    \textbf{Item} & \textbf{Reflection Bot M (SD)} & \textbf{Generic AI M (SD)} \\
    \midrule
    \endfirsthead
    \toprule
    \textbf{Item} & \textbf{Reflection Bot M (SD)} & \textbf{Generic AI M (SD)} \\
    \midrule
    \endhead
    \bottomrule
    \endfoot
    Überlegt, warum die eigenen Entscheidungen angemessen waren & 3.30 (0.90) & 3.90 (1.14) \\
    Mehrere Optionen verglichen, bevor entschieden wurde & 3.60 (0.92) & 3.80 (1.33) \\
    Erst während/danach klar geworden, warum die erste Idee nicht optimal war & 3.40 (0.66) & 2.50 (1.12) \\
    Konnte erklären, warum ein bestimmter Ansatz gewählt wurde & 3.60 (0.66) & 3.70 (1.10) \\
    System regte zu unabhängigerem Denken an als erwartet & 3.50 (1.20) & 2.40 (1.36) \\
    Hätte das jeweils andere System bevorzugt & 3.30 (1.27) & 2.20 (1.17) \\
    Arbeit an den Dilemma-Situationen hat die Fähigkeit verbessert, Prozessentscheidungen zu begründen & 3.10 (1.14) & 2.80 (1.25) \\
    War hilfreicher als erwartet & 3.40 (1.20) & 2.20 (1.33) \\
    Würde diese Art der KI-Unterstützung erneut nutzen & 2.50 (1.02) & 3.20 (1.89) \\
\end{longtable}
\endgroup

Die beiden zuletzt genannten Items bilden zusammen mit \glqq Kompetenzentwicklung\grqq{} und
\glqq anderes System bevorzugt\grqq{} die Grundlage von Abbildung~\ref{fig:survey-vergleich}.

\begingroup
\small
\setlength{\tabcolsep}{6pt}
\begin{longtable}{p{7.5cm} p{2.7cm} p{2.7cm}}
    \caption{Prä-Survey: Vergleichbarkeit der Gruppen zu Studienbeginn (n=9 Reflection Bot, n=7 Generic AI, nur eindeutig Prä-Post-zuordenbare Personen)}
    \label{tab:survey-praeexperience} \\
    \toprule
    \textbf{Item} & \textbf{Reflection Bot M (SD)} & \textbf{Generic AI M (SD)} \\
    \midrule
    \endfirsthead
    \toprule
    \textbf{Item} & \textbf{Reflection Bot M (SD)} & \textbf{Generic AI M (SD)} \\
    \midrule
    \endhead
    \bottomrule
    \endfoot
    Nutzt KI-Chatbots bereits für Lernaufgaben & 4.56 (0.83) & 4.14 (1.36) \\
    Weiß, wie man ein KI-System sinnvoll befragt & 4.33 (0.47) & 4.43 (0.49) \\
    Bewertet KI-Antworten kritisch, bevor sie übernommen werden & 4.44 (0.83) & 4.00 (0.76) \\
    Akzeptiert KI-Antworten meist direkt ohne Hinterfragen & 3.00 (1.33) & 2.57 (1.05) \\
\end{longtable}
\endgroup

Die beiden Gruppen sind hinsichtlich ihrer Vorerfahrung mit KI-Systemen bereits vor der Intervention
weitgehend vergleichbar. Die kleinen Unterschiede fallen dabei uneinheitlich aus: Die
Reflexionsbot-Gruppe bewertet KI-Antworten zwar etwas kritischer (4.44 gegenüber 4.00), berichtet
zugleich aber eine etwas höhere Tendenz, Antworten direkt zu übernehmen (3.00 gegenüber 2.57). Eine
eindeutig günstigere Ausgangslage der Reflexionsbot-Gruppe lässt sich daraus nicht ableiten. Die
Vorerfahrung erklärt spätere Gruppenunterschiede damit nicht zugunsten des Reflexionsbots.

\begin{table}[htbp]
    \centering
    \caption[Wissenstest zu Scrum und Prozessmodellen je Gruppe]{Wissenstest zu Scrum und
    Prozessmodellen je Gruppe (0 bis 6 Punkte, Prä-Post-zuordenbare Personen). Der
    Bewertungsschlüssel findet sich in \hyperref[app:fragebogen]{Anhang~A.5}.}
    \label{tab:survey-wissenstest}
    \begin{tabular}{p{3.5cm} p{2.8cm} p{2.8cm} p{1.2cm}}
        \toprule
        \textbf{Gruppe} & \textbf{Prä M (SD)} & \textbf{Post M (SD)} & \textbf{n} \\
        \midrule
        Reflection Bot & 5.56 (0.68) & 4.44 (1.89) & 9 \\
        Generic AI & 5.43 (0.73) & 5.43 (0.49) & 7 \\
        \bottomrule
    \end{tabular}
\end{table}

Etwa die Hälfte des Rückgangs in der Reflexionsbot-Gruppe entfällt auf einen einzelnen Fall, der
im Post-Survey alle sechs Fragen falsch beantwortete und in der qualitativen Kodierung unabhängig
davon bereits als Autorenschaft \glqq unklar\grqq{} geführt wird (vgl. \hyperref[app:kodierungstabelle]{Anhang~A.2}).
Die Gruppe verliert insgesamt zehn Punkte, fünf davon entfallen auf diesen Fall. Ohne ihn
liegt die Reflexionsbot-Gruppe bei Prä M=5.62, Post M=5.00 (n=8); ein leichter Rückgang bleibt damit
auch ohne ihn bestehen. Da beide Gruppen bereits im Prä-Test nahe der Höchstpunktzahl
liegen (Deckeneffekt), ist die Sensitivität des Wissenstests für Lernzuwächse ohnehin gering.

\begingroup
\small
\setlength{\tabcolsep}{5pt}
\begin{longtable}{p{5.3cm} p{1.7cm} p{1.7cm} p{1.7cm} p{1.7cm}}
    \caption{Reflexionsverhalten: Selbstauskunft Prä und Post je Gruppe (n=9 Reflection Bot, n=7 Generic AI)}
    \label{tab:survey-reflexionsverhalten} \\
    \toprule
    \textbf{Item} & \textbf{RB Prä} & \textbf{RB Post} & \textbf{AI Prä} & \textbf{AI Post} \\
    \midrule
    \endfirsthead
    \toprule
    \textbf{Item} & \textbf{RB Prä} & \textbf{RB Post} & \textbf{AI Prä} & \textbf{AI Post} \\
    \midrule
    \endhead
    \bottomrule
    \endfoot
    Überlegt, warum Entscheidung angemessen war & 3.67 & 3.44 & 3.57 & 4.00 \\
    Mehrere Optionen verglichen vor Entscheidung & 4.11 & 3.56 & 4.43 & 3.86 \\
    Erst im Nachhinein klar, warum erste Idee suboptimal war & 3.78 & 3.44 & 2.86 & 2.29 \\
    Konnte eigenen Ansatz begründen & 4.00 & 3.56 & 4.43 & 3.71 \\
    Reflektierte über Feedback bzw. Gegenfragen & 4.00 & 3.22 & 4.43 & 3.43 \\
\end{longtable}
\endgroup

Die Reflexionsbot-Gruppe schätzt ihr Reflexionsverhalten bei allen fünf Items nach der Übung niedriger
ein als vorher; in der Generic-AI-Gruppe gilt das für vier von fünf Items ebenso. Eine mögliche
Erklärung ist ein Response-Shift-Effekt (Antwortverschiebungseffekt): Die Selbsteinschätzung vor der Aufgabe war eher
abstrakt optimistisch, nach der konkreten Konfrontation mit den Dilemma-Situationen wird der eigene
Anspruch realistischer kalibriert. Ein eigenständiger Effekt des Reflexionsbots auf das
selbstberichtete Reflexionsverhalten ist auf dieser Basis nicht erkennbar. Der in
Abschnitt~\ref{sec:beantwortung-forschungsfragen} berichtete Unterschied beim \glqq unabhängigeren
Denken als erwartet\grqq{} bezieht sich auf eine andere, erwartungsbezogene Frage und steht dazu nicht
im Widerspruch.

\begingroup
\small
\setlength{\tabcolsep}{6pt}
\begin{longtable}{p{8.5cm} p{2cm} p{2cm}}
    \caption{Selbsteingeschätztes Prozessverständnis: Veränderung Post minus Prä je Gruppe (Skala 1 bis 5)}
    \label{tab:survey-prozessverstaendnis} \\
    \toprule
    \textbf{Stufe / Item} & \textbf{RB $\Delta$ (n=9)} & \textbf{AI $\Delta$ (n=7)} \\
    \midrule
    \endfirsthead
    \toprule
    \textbf{Stufe / Item} & \textbf{RB $\Delta$ (n=9)} & \textbf{AI $\Delta$ (n=7)} \\
    \midrule
    \endhead
    \bottomrule
    \endfoot
    Stufe 1: Prozessmodell erklären & $-$0.56 & $-$0.71 \\
    Stufe 1: Sprint-Ablauf beschreiben & $-$0.33 & $-$0.71 \\
    Stufe 1: Unterschiede agil/plan-basiert benennen & $-$0.78 & $-$0.86 \\
    Stufe 2: Eignung eines Prozessmodells begründen & $+$0.33 & $-$0.71 \\
    Stufe 2: kurze Sprints begründen & $+$0.11 & $-$0.43 \\
    Stufe 2: Schnelligkeit vs. Dokumentation erklären & 0.00 & $-$0.57 \\
    Stufe 3: Prozessmodell-Wahl abwägen & $+$0.11 & $-$0.86 \\
    Stufe 3: Anpassungsbedarf bei rechtlichen Vorgaben erkennen & $+$0.11 & $-$0.57 \\
\end{longtable}
\endgroup

Bei den Stufe-1-Items (reines Beschreiben) verlieren beide Gruppen leicht, was zum oben beschriebenen
Response-Shift-Effekt passt. Bei den Stufe-2- und Stufe-3-Items (Begründen, Abwägen) sinkt die
Selbsteinschätzung in der Generic-AI-Gruppe dagegen durchgehend weiter, während sie in der
Reflexionsbot-Gruppe stabil bleibt oder leicht steigt, wie in
Abschnitt~\ref{sec:beantwortung-forschungsfragen} berichtet.

\begingroup
\small
\setlength{\tabcolsep}{6pt}
\begin{longtable}{p{7.5cm} p{2.7cm} p{2.7cm}}
    \caption{Post-Survey: Kognitive Last (Skala 1 = very easy / very low bis 5 = very difficult /
    very high; n=10 Reflection Bot, n=10 Generic AI)}
    \label{tab:survey-kognitive-last} \\
    \toprule
    \textbf{Item} & \textbf{Reflection Bot M (SD)} & \textbf{Generic AI M (SD)} \\
    \midrule
    \endfirsthead
    \toprule
    \textbf{Item} & \textbf{Reflection Bot M (SD)} & \textbf{Generic AI M (SD)} \\
    \midrule
    \endhead
    \bottomrule
    \endfoot
    Task 1: Analyse der Scrum-Elemente & 3.00 (1.00) & 3.00 (1.55) \\
    Task 2: Entscheidung in den Dilemma-Situationen & 2.70 (1.00) & 2.50 (1.20) \\
    Task 3: Reflexion der schwierigsten Entscheidung & 2.90 (1.04) & 2.50 (1.12) \\
    Das KI-System effektiv nutzen & 2.80 (1.17) & 2.70 (1.62) \\
    Die Gegenargumente des KI-Systems bewerten & 3.10 (1.22) & 2.80 (1.66) \\
    \midrule
    \textbf{Gesamter kognitiver Aufwand der Übung} & \textbf{3.00 (1.10)} & \textbf{3.20 (1.08)} \\
\end{longtable}
\endgroup

Bei den fünf Einzelitems liegt die Reflexionsbot-Gruppe zwischen 0.00 und 0.40 Skalenpunkten höher,
bei der zusammenfassenden Einschätzung des gesamten kognitiven Aufwands dagegen um 0.20 Punkte
niedriger. Sämtliche Unterschiede sind kleiner als die jeweilige Standardabweichung. Die Skala stützt
damit weder die Annahme, der Reflexionsbot sei als anstrengender empfunden worden, noch deren
Gegenteil (vgl. Abschnitt~\ref{sec:ergebnisse}).

\newpage
\section*{A.7 Systemnachricht des Reflexionsbots}
\phantomsection
\label{app:systemprompt}
\addcontentsline{toc}{section}{A.7 Systemnachricht des Reflexionsbots}

Die folgende Systemnachricht (Systemprompt) wurde dem Sprachmodell \texttt{gpt-4o-mini} bei jeder
Anfrage über die Rolle \texttt{system} übergeben (vgl. Abschnitt~\ref{sec:technische-umsetzung}) und
setzt die in Abschnitt~\ref{sec:konzept-reflexionsbot} beschriebenen didaktischen Prinzipien
(sokratisches Fragen, adaptive Tiefe, Szenariobezug ohne Namensnennung) technisch um. Sie ist hier im
englischen Original wiedergegeben. Einzelne Unicode-Sonderzeichen (Pfeile, Gedankenstriche) sind aus
Satzgründen als ASCII-Zeichen dargestellt.

\begin{lstlisting}[numbers=none, breaklines=true, basicstyle=\scriptsize\ttfamily]
You are a didactic AI reflection coach for software engineering students working on process modelling scenarios.

CONTEXT -- TWO SCENARIOS THE STUDENTS ARE WORKING ON:

SCENARIO 1 -- SMART INSULIN PEN (SmartDose GmbH)
A company adds an ML dosage-recommendation model to a connected insulin pen for diabetic patients.
Key facts students must grapple with:
- 9-person Scrum team, 2-week sprints, no ML or data science experience on the team
- ML output is probabilistic -- wrong insulin dose is life-threatening
- EU Medical Device Regulation (MDR) requires clinical validation before any real patient sees the output
- Training data requires sensitive personal health data (blood glucose, meal logs, activity)
- CEO wants a patient-facing prototype in 6 months; MDR approval takes minimum 18 months
- Key decisions: which 2 of 4 roles to hire first (ML Engineer, Clinical Validator, Data Privacy Officer, Regulatory Specialist); which dataset to use (50k records without consent vs 5k records with consent); how to respond to CEO's impossible 6-month timeline

SCENARIO 2 -- LEARNLOOP (Distributed Startup)
An AI-powered adaptive learning platform startup with teams in Amsterdam (UTC+1), Bangalore (UTC+5:30), and Austin (UTC-6).
Key facts:
- 12 people, 9-month MVP deadline, 3.2M EUR raised
- Bangalore-Austin overlap is nearly zero (11.5h gap); Amsterdam-Bangalore overlap is small
- University stakeholders (IT, Professors, Students, Legal) disagree on requirements
- Austin collects requirements weekly but has no process to pass them to the other teams
- Bangalore built a recommendation engine before Amsterdam decided what data to collect
- Key decisions: Amsterdam vs Bangalore architecture conflict (whose design wins); which agenda item to drop from the only 90-min all-team slot; what to do when a FERPA legal blocker freezes all student data collection mid-sprint

YOUR ROLE AND RULES:
- Always respond in English, regardless of which language the student writes in. Students may write in German or English -- both are fine. Never switch your response language to German.
- You guide students using Socratic scaffolding -- ask targeted questions, give short hints, NEVER give the full answer.
- ADAPTIVE DEPTH -- this is your most important rule. The richness of your response must scale directly with the quality of the student's reasoning:

    LEVEL 1 -- Vague or minimal answer (e.g. "I would pick option 2", "maybe hire the ML engineer"):
    -> Give ONLY 1 short question to pull out their reasoning. Nothing more.
       Example: "What specific risk are you trying to avoid with that choice?"

    LEVEL 2 -- Partial reasoning (student names a choice and gives one reason, but skips risks or constraints):
    -> Briefly confirm what they got right (1 sentence), then ask 1-2 focused questions about what they missed.
       Example: "That covers the technical side well -- but what happens to the project if that role is missing at the point where approval needs to start?"

    LEVEL 3 -- Strong reasoning (student mentions trade-offs, acknowledges a real constraint, considers what they lose with their choice, thinks about timing or stakeholders):
    -> REWARD them visibly: start by validating the strongest point they made (be specific, not generic praise).
       Then add one insight or perspective they did not mention -- something that genuinely extends their thinking.
       Close with 1-2 higher-order questions that push them to the next level.
       This response can be longer and more substantive -- the student earned it.

    LEVEL 4 -- Exceptional answer (student argues a trade-off from multiple angles, references a real constraint by name or detail, anticipates a second-order effect):
    -> Full engagement: confirm what is genuinely strong, surface the single best counterargument to their position, and ask one question that challenges their core assumption.
       Treat them as a peer, not a student being corrected.
- Pick ONE logical weakness per round to focus on (e.g., first round: risk, next round: roles/stakeholders, then dependencies, then timeline realism).
- Keep responses concise -- no walls of text even for strong answers. Three well-aimed sentences beat ten generic ones.
- If the student pastes in the AI Analysis prompt result (from ChatGPT/Claude/Gemini), help them critically evaluate whether the AI counterargument was valid or ignored their constraints.
- NEVER explicitly mention the scenario names or numbers ("Scenario 1", "Scenario 2", "SmartDose", "LearnLoop") in your responses. Use your knowledge of them silently to ask relevant questions and give accurate feedback -- but the student should not be able to tell which scenario you know about from your response.
- If input is completely off-topic or unclear, politely ask them to clarify what they are working on.
- Use a friendly, clear, academically rigorous tone. Treat students as capable thinkers who need a sharper mirror, not a lecture.
\end{lstlisting}

\newpage
\section*{A.8 Verwendete KI-Werkzeuge}
\phantomsection
\label{app:kiwerkzeuge}
\addcontentsline{toc}{section}{A.8 Verwendete KI-Werkzeuge}

Gemäß der eidesstattlichen Erklärung, der Prüfungsordnung der Hochschule~Heilbronn sowie den
Leitlinien guter wissenschaftlicher Praxis \cite{dfg2022codeofconduct, verch2024prompt} wird der
Einsatz \gls{ki}-gestützter Werkzeuge während der Erstellung dieser Arbeit nach Zweck und Umfang in
Tab.~\ref{tab:ki-werkzeuge} dokumentiert.
Alle damit erzeugten Inhalte wurden vom Autor geprüft, verifiziert und inhaltlich überarbeitet.
Kein KI-Werkzeug wurde als ungeprüfte Quelle für Sachaussagen verwendet.

\begingroup
\small
\setlength{\tabcolsep}{6pt}
\renewcommand{\arraystretch}{1.15}
\begin{longtable}{%
        >{\raggedright\arraybackslash}p{3.2cm}%
        >{\raggedright\arraybackslash}p{3.9cm}%
        >{\raggedright\arraybackslash}p{7.0cm}}
    \caption{Verwendete KI-Werkzeuge. Alle Ausgaben wurden vom Autor geprüft und validiert.}
    \label{tab:ki-werkzeuge} \\
    \toprule
    \textbf{Werkzeug} & \textbf{Einsatzzweck} & \textbf{Umfang der Nutzung} \\
    \midrule
    \endfirsthead
    \toprule
    \textbf{Werkzeug} & \textbf{Einsatzzweck} & \textbf{Umfang der Nutzung} \\
    \midrule
    \endhead
    \bottomrule
    \endfoot

    \textbf{Claude Code}\newline
    {\footnotesize Claude Sonnet 4.6 und 5, Claude Opus 4.7, 4.8 und 5.0} &
    Code, Experimente, Text und Abbildungen, Beratung &
    Hauptassistent ab den Umsetzungsphasen an. Unterstützte den Prozess des Verfassens und 
    der Fehlerbehebung bei den meisten Codestellen und Code-Anpassungen für den Reflexionsbot, diskutierte das Design vom codebasierten Prototypen
    und andere Entscheidungen, interpretierte Ergebnisse, half 
    beim Verfassen und Überarbeiten der Kapitel der Bachelorarbeit 
    unter Berücksichtigung der Schreibrichtlinien, half bei der Erstellung der Diagramme 
    und Abbildungen sowie bei weiteren unterstützenden Aufgaben. \\
    \midrule

    \textbf{GitHub Copilot}\newline
    {\footnotesize Claude Sonnet 4.6, automatische Modellauswahl} &
    Code- und Tooling-Setup, Recherche\-unterstützung &
    Punktuell in der Entwicklungsphase; Verbesserung der Qualität und Effizienz des Python-Codes; Beschaffung zusätzlicher Informationen; Hilfe beim Schreiben und Optimieren für die explorative Datenanalyse\\
    \midrule

    \textbf{ResearchRabbit} &
    Literaturrecherche und Nachverfolgung &
    Erweiterung des Quellenbestands über die Zitationsbeziehungen
    bereits gefundener Veröffentlichungen. \\
    \midrule

    \textbf{NotebookLM}\newline
    {\footnotesize Gemini 3} &
    Literatur\-organisation und Recherche\-unterstützung &
    Zusammenfassen und Strukturieren der Veröffentlichungen sowie
    Unterstützung bei Verständnisfragen. \\
\end{longtable}
\endgroup

\section*{A.9 Digitaler Anhang}
\phantomsection
\label{app:digitaleranhang}
\addcontentsline{toc}{section}{A.9 Digitaler Anhang}

Zur Nachvollziehbarkeit der Untersuchung sind der Quellcode des Reflexionsbots, die
Erhebungsmaterialien und die \LaTeX-Quellen dieser Arbeit im begleitenden Repository abgelegt:

\begin{itemize}
    \item Persönliches GitHub (öffentlich):\\
    \url{https://github.com/Bastipreme/Bachelorarbeit}

    \item Hochschule Heilbronn GitLab (institutionell):\\
    \url{https://git.it.hs-heilbronn.de/seschuster/bachelorarbeit}
\end{itemize}

Die Rohdaten des Mini-Experiments sind aus Gründen des Teilnehmendenschutzes nicht Bestandteil des
öffentlichen Repositorys (vgl. Abschnitt~\ref{sec:forschungsethik}). Repository und digitaler Anhang
sind wie folgt gegliedert:

\begin{itemize}
    \item \texttt{Chatbot/} enthält die Implementierung des Reflexionsbots. \texttt{Chatbot.py}
    umfasst die Streamlit-Anwendung mit Systemnachricht, Sitzungsverwaltung und Anbindung an die
    Chat-Completions-API (vgl. Abschnitt~\ref{sec:technische-umsetzung}), \texttt{Dockerfile} und
    \texttt{docker-compose.yml} den Container, über den die Anwendung in der Übung erreichbar war,
    und \texttt{requirements.txt} die verwendeten Paketversionen. Der OpenAI-API-Schlüssel ist nicht
    enthalten. An seiner Stelle liegen die Vorlagen \texttt{secrets.toml.example} und
    \texttt{.env.example} bei, die ohne die Endung \texttt{.example} zu kopieren und mit einem
    eigenen Schlüssel zu füllen sind. Das Verhalten des Bots ändert sich dadurch nicht, da Modell
    und Systemnachricht im Code festgelegt sind (vgl. \hyperref[app:systemprompt]{Anhang~A.7}).

    \item \texttt{Experiment/} enthält die Erhebungsmaterialien: die Aufgabenblätter beider
    Szenarien in je einer Fassung für die Reflexionsbot- und die Generic-AI-Gruppe (vgl.
    \hyperref[app:szenarien]{Anhang~A.4}), die Fragebogentexte für Prä- und Post-Erhebung
    (vgl. \hyperref[app:fragebogen]{Anhang~A.5}) sowie die Skripte, mit denen die zugehörigen
    Online-Formulare erzeugt wurden.

    \item \texttt{Thesis/} enthält die \LaTeX-Quellen dieses Dokuments einschließlich der
    Literaturdatenbank und der eingebundenen Abbildungen.

    \item \texttt{Data/} liegt ausschließlich im digitalen Anhang der eingereichten Fassung, nicht
    im öffentlichen Repository. Es enthält die beiden Tabellendateien mit den vollständigen
    Antworten des Prä- und des Post-Surveys sowie den Ordner \texttt{Upload-Assignments/} mit den
    eingereichten Bearbeitungen und Chatverläufen, getrennt nach Gruppen und je Person
    pseudonymisiert (vgl. \hyperref[app:kodierungstabelle]{Anhang~A.2}).
\end{itemize}
